\documentclass[10pt, conference]{IEEEtran}
\usepackage{cite}
\usepackage{amsmath,amssymb,amsfonts}
\usepackage{algorithmic}
\usepackage{graphicx}
\usepackage{textcomp}
\usepackage{xcolor}
\usepackage{hyperref}

\begin{document}

\title{Bridging the Performance Gap: Upgrading Deep SAD with End-to-End Spectrogram CNNs for Robust Zero-Day Attack Detection in TF-QKD}

\author{\IEEEauthorblockN{1\textsuperscript{st} Anonymous Author}
\IEEEauthorblockA{\textit{Department of Physics and Computer Science} \\
\textit{Institute of Quantum Cryptography}\\
City, Country \\
email@domain.com}
}

\maketitle

\begin{abstract}
Recent statistical evaluations of Twin-Field Quantum Key Distribution (TF-QKD) physical-layer security have revealed a significant performance disparity between tree-ensemble models and Deep Semi-Supervised Anomaly Detection (Deep SAD). While XGBoost achieves a 97.4\% AUROC on zero-day attack holdouts, Deep SAD trails due to representation overfitting and hypersphere collapse, culminating in a 0\% True Positive Rate against the Carrier Density Attack (CDA). This report documents a comprehensive architectural overhaul designed to enable Deep SAD to surpass XGBoost. By replacing the 10-dimensional manual feature extraction bottleneck with an end-to-end 2D Convolutional Neural Network (CNN) trained directly on Short-Time Fourier Transform (STFT) spectrograms, the model is empowered to discover subtle amplitude perturbations autonomously. Furthermore, we document the integration of Robust Scaling and a Margin-Based Contrastive Loss to strictly enforce the hypersphere decision boundary.
\end{abstract}

\begin{IEEEkeywords}
TF-QKD, Deep SAD, Anomaly Detection, Spectrogram, CNN, Carrier Density Attack
\end{IEEEkeywords}

\section{Introduction}
The security of Twin-Field Quantum Key Distribution (TF-QKD) against Trojan-horse and optical injection-locking side-channels relies heavily on the robust detection of physical-layer anomalies. An extensive empirical evaluation of our intrusion detection system revealed that supervised baselines, specifically XGBoost, significantly outperformed the semi-supervised Deep SAD framework in identifying zero-day attacks. 

Feature importance analysis (via SHapley Additive exPlanations and Permutation Drop) illuminated the root cause of this disparity: the models anchor their decision boundaries heavily on phase-based statistics, predominantly \textit{phase kurtosis} and \textit{lag-1 phase autocorrelation}. Consequently, attacks that exclusively manipulate amplitude dynamics, such as the Carrier Density Attack (CDA), entirely evade detection, resulting in a 0\% True Positive Rate (TPR) at a 4\% False Positive Rate (FPR) threshold for both models. 

To resolve this blindspot and elevate Deep SAD beyond the capabilities of XGBoost, we propose an architectural paradigm shift. XGBoost is intrinsically limited by the manual 10-dimensional statistical features provided to it, as tree ensembles cannot effectively process raw sequence data of 100,000 pulses. Conversely, Deep Learning architectures excel at high-dimensional sequence modeling. This document formally outlines the implementation plan to upgrade Deep SAD from a rudimentary Multi-Layer Perceptron (MLP) to an end-to-end Spectrogram-based 2D-CNN, alongside critical mathematical refinements to the loss objective and data normalization pipeline.

\section{Proposed Methodology}

The proposed upgrade consists of four structural modifications to the Deep SAD pipeline.

\subsection{Time-Frequency Feature Representation}
\textbf{Objective:} Transition from manual statistical features to rich, uncompressed physical representations.

\textbf{Methodology:} The current pipeline reduces windows of 100,000 raw physical pulses into a 10-dimensional scalar vector. We will introduce an `extract_spectrogram` function utilizing the Short-Time Fourier Transform (STFT) via SciPy. This will process the raw intensity traces ($I_A$, $I_B$) and phase difference ($\Delta\phi$), producing a 3-channel time-frequency grid (Channels $\times$ Frequencies $\times$ Time Steps).

\textbf{Alternative Considered:} Feeding the raw 100,000-point 1D time-series into a 1D-CNN (e.g., WaveNet) was considered. However, the immense sequence length renders standard 1D convolutions computationally prohibitive without aggressive pooling, which risks destroying high-frequency physical phenomena. Spectrograms naturally compress the temporal domain while preserving the spectral richness necessary to detect anomalies like CDA.

\subsection{Spectrogram-based 2D-CNN Architecture}
\textbf{Objective:} Enable autonomous feature discovery of high-frequency amplitude and phase perturbations.

\textbf{Methodology:} The `DeepSVDD` class will be restructured to utilize a `SpectrogramCNN` constructed in PyTorch. The architecture will consist of three 2D Convolutional layers, each followed by Batch Normalization, LeakyReLU activations, and MaxPooling. The extracted feature maps will be flattened and linearly projected into the final 8-dimensional hypersphere representation space. 

\textbf{Alternative Considered:} Vision Transformers (ViT) and Long Short-Term Memory (LSTM) networks were evaluated. LSTMs are computationally infeasible for this sequence length. ViTs lack the inductive bias of translation invariance and are prone to severe representation collapse in semi-supervised anomaly detection without exhaustive pre-training on massive datasets. The CNN offers the optimal balance of inductive bias and training efficiency.

\subsection{Robust Scaling for Outlier Mitigation}
\textbf{Objective:} Prevent the artificial collapse of the nominal hypersphere caused by extreme feature outliers.

\textbf{Methodology:} The existing implementation utilizes `StandardScaler`, which standardizes features to zero mean and unit variance. Because advanced injection attacks induce massive spikes in parameters like phase kurtosis, `StandardScaler` severely compresses the variance of the \textit{nominal} data, forcing the neural network to collapse the hypersphere boundary artificially. This will be replaced with `RobustScaler`, which relies on the median and interquartile range (IQR), rendering the normalization immune to attack-induced outliers.

\subsection{Margin-Based Semi-Supervised Contrastive Loss}
\textbf{Objective:} Enforce a strict, mathematically bounded decision boundary to separate nominal and anomalous distributions.

\textbf{Methodology:} The original Deep SAD loss relies on an inverse distance penalty for anomalous samples: $\mathcal{L}_{anom} = \eta / (d + \epsilon)$, where $d$ is the squared Euclidean distance to the hypersphere center. This unbounded penalty yields vanishing gradients at large distances and produces fuzzy decision boundaries, contributing to the low TPR. We will implement a Margin-Based Contrastive Loss:
\begin{equation}
\mathcal{L}_{anom} = \max(0, m - d)
\end{equation}
where $m$ is a tunable margin hyperparameter. This formulation strictly forces known anomalies outside the defined hypersphere radius, creating a rigid mathematical "wall."

\section{Experimental Pipeline and Integration}

To accurately evaluate the upgraded Deep SAD against the classical baseline, the data loading logic in `models/train.py` will be bifurcated. XGBoost will continue to receive the 10-dimensional statistical vectors derived from `extract_features`, as tree-based models cannot natively process 3D image arrays. Simultaneously, Deep SAD will be fed the 3-channel physical spectrograms. 

This parallel routing ensures that both architectures are evaluated on their optimal input representations, allowing us to definitively test the hypothesis that end-to-end deep learning on raw physical waveforms can supersede manually engineered statistical features in detecting zero-day physical-layer attacks.

\section{Conclusion}
The modifications outlined in this report address the fundamental algorithmic limitations of the current Deep SAD implementation. By upgrading to a Spectrogram-CNN with margin-based contrastive loss and robust scaling, we anticipate a significant increase in True Positive Rates for previously evasive amplitude-based attacks, positioning Deep SAD as the superior detection mechanism for TF-QKD security monitoring.

\end{document}
